\subsection{Affine growth and projective envelopes}
\label{sec:affine-projective-envelopes}
\index{Envelope!affine growth}
\index{Contact set!growth parameter}

An obstacle imposed on a compact subset of an affine chart can also be studied on projective space. The logarithmic growth becomes part of the background. This correspondence separates contact with the raw supremum from contact with its upper semicontinuous regularization.

In this subsection write $\chi(z)=\log(1+\lVert z\rVert^2)$ on $\mathbb C^n$, and let $\omega_{\mathrm{FS}}$ be the Fubini--Study form on $\mathbb P^n$ with this local potential. Thus $\int_{\mathbb P^n}\omega_{\mathrm{FS}}^n=1$ in the book convention. For $\delta>0$ put
\[
 \theta_\delta=\tfrac\delta2\omega_{\mathrm{FS}},\qquad
 \mathcal L_\delta=
 \{u\in\PSH(\mathbb C^n):
    u(z)\leq\delta\log^+\lVert z\rVert+C_u\}.
\]
Here $\log^+0=0$. Let $\mathcal L_\delta^+$ consist of the members also bounded below by $\delta\log^+\lVert z\rVert-C$ for some constant $C$.

\Needspace{8\baselineskip}
\begin{lemma}\label{lem:affine-projective-growth}
The map $u\mapsto u-\delta\chi/2$, followed by canonical extension across the hyperplane at infinity, identifies $\mathcal L_\delta$ with $\PSH(\mathbb P^n,\theta_\delta)$.
\end{lemma}
\begin{proof}
Use homogeneous coordinates $[Z_0:\cdots:Z_n]$ and the affine chart $Z_0\ne0$. On a chart $Z_i\ne0$ meeting $H=\{Z_0=0\}$, the local plurisubharmonic representative of $p=u-\delta\chi/2$ is
\[
 p+\frac\delta2\log\frac{\sum_{j=0}^n|Z_j|^2}{|Z_i|^2}
 =u+\delta\log|Z_0/Z_i|.
\]
It is locally bounded above near $H$: the norm of the affine coordinates is at most a fixed chart constant divided by $|Z_0/Z_i|$. The growth bound and \cref{lem:qpsh-removable-pluripolar} therefore give a canonical extension. The smooth potential transitions make the extensions agree on overlaps.

Conversely, a compact $\theta_\delta$-plurisubharmonic $p$ is bounded above. The function $u=p+\delta\chi/2$ is plurisubharmonic on the affine chart, and
\[
 0\leq\tfrac12\chi-\log^+\lVert z\rVert\leq\tfrac12\log2
\]
gives the required upper growth. Neither restriction is identically $-\infty$. Canonical extension gives uniqueness in both directions.
\end{proof}

Fix a non-pluripolar compact $K\subset\mathbb C^n$ and a bounded real function $\psi$ on $K$. Define
\[
 V_\delta=\sup\{u\in\mathcal L_\delta:u|_K\leq\psi\},
 \qquad F_\delta=V_\delta^*,
\]
where the obstacle is required everywhere on $K$ before regularization. On $\mathbb P^n$ put
\[
 G_\delta=\sup\{p\in\PSH(\mathbb P^n,\theta_\delta):
                 p|_K\leq\psi-\delta\chi/2\}.
\]

\Needspace{10\baselineskip}
\begin{proposition}\label{prop:affine-projective-envelope}
The compact envelope $G_\delta^*$ is bounded and $\theta_\delta$-plurisubharmonic. On the affine chart,
\begin{equation}\label{eq:affine-projective-envelope}
 V_\delta=\tfrac\delta2\chi+G_\delta,\qquad
 F_\delta=\tfrac\delta2\chi+G_\delta^*\in\mathcal L_\delta^+.
\end{equation}
The raw envelope agrees with its regularization almost everywhere.
\end{proposition}
\begin{proof}
The lemma identifies the competitor families and proves the raw identity. Constants below $\inf_K(\psi-\delta\chi/2)$ supply compact competitors.

The constant upper obstacle $\sup_K(\psi-\delta\chi/2)$ and \cref{lem:compact-obstacle-upper-bound} give a common upper bound for the competitors. A constant lower competitor and the upper-envelope theorem then give bounded $G_\delta^*$. The smooth addition in \eqref{eq:affine-projective-envelope} commutes with regularization and gives the two-sided growth bounds.

Locally adding a potential of $\theta_\delta$ makes the competitor family plurisubharmonic. The upper-envelope theorem recalled in \cref{chap:framework} says that regularization changes its supremum only on a pluripolar set. Thus $G_\delta=G_\delta^*$ almost everywhere; smooth addition gives the same assertion for $V_\delta$.
\end{proof}

Pluripolar exceptions can also be removed within the logarithmic growth class. A pole barrier with the same growth allowance gives a direct repair of each competitor.

\begin{lemma}\label{lem:affine-logarithmic-pole-barrier}
For every pluripolar $E\subseteq\mathbb C^n$ and every $\delta>0$, there is a genuine $b\in\mathcal L_\delta$ with $b=-\infty$ on $E$. Given a non-pluripolar compact $K$, it can be chosen with $b\leq0$ on $K$.
\end{lemma}
\begin{proof}
Put $E_j=E\cap\overline B(0,j)$. Each $E_j$ is pluripolar on $\mathbb P^n$: it has local witnesses at every affine point, while a neighborhood of the hyperplane at infinity avoids this bounded set. By the compact global Josefson theorem, \cite[Theorem~1.1.5]{SGPT}, there is a genuine quasi-plurisubharmonic pole function for $E_j$. Rescaling its curvature bound and subtracting its supremum give
\[
 p_j\in\PSH(\mathbb P^n,\theta_1),\qquad
 p_j\leq0,\qquad p_j|_{E_j}=-\infty.
\]
The $L^1$ norms of these normalized potentials are uniformly bounded by \cref{thm:normalized-compactness}. The sums
\[
 p_N=\sum_{j=1}^N2^{-j}p_j
\]
decrease, have uniformly bounded $L^1$ norms, and remain $\theta_1$-plurisubharmonic because $\theta_1\geq0$ and the coefficients sum to at most one. \Cref{cor:decreasing-limit} gives a genuine canonical limit $p$. Every point of $E$ is a pole of one summand, hence $p=-\infty$ there. By \cref{lem:affine-projective-growth},
\[
 b=\delta(p+\chi/2)\in\mathcal L_\delta
\]
is the required affine pole function. Its supremum on $K$ is finite: compactness bounds it above, and non-pluripolarity excludes $b=-\infty$ on all of $K$. Subtracting this supremum gives $b\leq0$ on $K$.
\end{proof}
\index{Pluripolar set!logarithmic pole barrier}

\begin{corollary}\label{cor:affine-polar-obstacle-invariance}
The envelope $F_\delta$ is unchanged if $u|_K\leq\psi$ is imposed only quasi-everywhere, or only on $K\setminus E$ for any pluripolar $E\subseteq K$.
\end{corollary}
\begin{proof}
Let $u\in\mathcal L_\delta$ satisfy $u\leq\psi$ on $K\setminus E$, where $E$ is pluripolar. Choose $b$ from the lemma with this pole set and $b\leq0$ on $K$, and put $m=\inf_K\psi$. For $\epsilon>0$ define
\[
 u_\epsilon=\frac{u+\epsilon b+\epsilon m}{1+\epsilon}.
\]
This is a genuine member of $\mathcal L_\delta$. On $K\setminus E$, the inequalities $u\leq\psi$, $b\leq0$ and $m\leq\psi$ give $u_\epsilon\leq\psi$; on $E$ the pole term enforces the constraint. Thus $u_\epsilon\leq F_\delta$ and $u_\epsilon\to u$ almost everywhere. Canonical order gives $u\leq F_\delta$ everywhere. Taking the supremum and regularization proves the asserted equalities.

Equivalently, apply \cref{prop:compact-obstacle-polar-invariance} on $\mathbb P^n$ with the big background $\theta_\delta$ and the bounded obstacle $\psi-\delta\chi/2$, then use \eqref{eq:affine-projective-envelope}.
\end{proof}

These equalities concern regularizations. The raw envelope $V_\delta$ continues to use the obstacle at every point of $K$.


Changing $\delta$ changes both $\theta_\delta$ and the obstacle $\psi-\delta\chi/2$. The raw contact condition nevertheless has a simple parameter structure. Put
\[
 \Lambda=\{\delta>0:V_\delta|_K=\psi\},\qquad
 \Lambda_{\mathrm{reg}}(a)=
 \{\delta\in\Lambda:F_\delta(a)=\psi(a)\}\quad(a\in K).
\]

\Needspace{10\baselineskip}
\begin{theorem}\label{thm:affine-growth-contact-ray}
The contact set of parameters $\Lambda$ is empty, $(0,\infty)$, or $[\lambda,\infty)$ for some $\lambda>0$. If $\Lambda=[\lambda,\infty)$, each $\Lambda_{\mathrm{reg}}(a)$ is empty, $\{\lambda\}$, or $\Lambda$. If $\Lambda=(0,\infty)$, it is empty or $\Lambda$. The same alternatives hold for $\bigcap_{a\in K}\Lambda_{\mathrm{reg}}(a)$.
\end{theorem}
\begin{proof}
Let $m=\inf_K\psi$. Inclusion of growth classes and scaling of competitors give, for $0<\alpha<\beta$,
\begin{equation}\label{eq:affine-growth-raw-slope}
 0\leq V_\beta-V_\alpha
 \leq\frac{\beta-\alpha}{\alpha}(V_\alpha-m).
\end{equation}
Indeed, if $u$ is a $\beta$-competitor, then $(\alpha/\beta)(u-m)+m$ is an $\alpha$-competitor, since $\psi\geq m$. The constant $m$ is also a competitor, so all terms are finite and the right side is nonnegative. The estimate gives continuity of each raw value in the positive parameter. Since $V_\delta\leq\psi$ on $K$, monotonicity makes $\Lambda$ upward closed; continuity makes it relatively closed in $(0,\infty)$. This proves its three possible forms.

For $0\leq t\leq1$, convex combinations of competitors give
\[
 V_{(1-t)\alpha+t\beta}\geq(1-t)V_\alpha+tV_\beta.
\]
The preceding proposition turns this into the same inequality for the regularized envelopes almost everywhere, then everywhere by \cref{lem:qpsh-almost-everywhere-order}. Thus $F_\delta$ is concave and nondecreasing in $\delta$ at every point.

For $a\in K$ and $\delta\in\Lambda$, the defect $d_\delta=F_\delta(a)-\psi(a)$ is nonnegative, concave and nondecreasing. Suppose $d_\alpha=0$ and choose $\eta\in\Lambda$ with $\eta<\alpha$. For every $\beta>\alpha$, express $\alpha$ as a strict convex combination of $\eta$ and $\beta$. Concavity forces both nonnegative endpoint defects to vanish. Monotonicity gives vanishing at all smaller parameters in $\Lambda$. Hence any zero other than the smallest contact parameter forces every defect to vanish. This proves the pointwise alternatives. Their intersection gives the asserted global alternatives.
\end{proof}

The distinction concerns the regularized value at a point of $K$. Raw contact persists as the growth allowance increases; regularized contact can be restricted to the first positive contact parameter. The compact correspondence keeps the background and obstacle matched throughout this variation.

\subsubsection{Growth slopes and supporting probabilities}
\label{sec:affine-envelope-slopes}

By concavity, the rate at which $F_\delta$ grows when the growth allowance $\delta$ is increased is at least its rate of growth for large $\delta$, and the latter is the classical extremal function of $K$, without weight. For continuous weights, supporting probabilities express the initial rate of growth through tests by functions of the fixed class $\mathcal L_1$.

When the weight varies, we write $V_{\delta,q}$ and $F_{\delta,q}=V_{\delta,q}^*$ for the raw and regularized envelopes defined as above, with a bounded obstacle $q$ on the fixed compact set $K$ in place of $\psi$. Thus $F_\delta=F_{\delta,\psi}$. Set
\[
 V_K=V_{1,0},\qquad V_K^*=F_{1,0},\qquad
 m=\inf_K\psi,\qquad M=\sup_K\psi.
\]
Thus $V_K^*$ is the extremal function of $K$. Since the constant $0$ is a competitor for $V_K$, and since competitors can be scaled and translated by constants, we have $V_K^*\geq0$ and
\begin{equation}\label{eq:affine-obstacle-sandwich}
 \delta V_K^*+m\leq F_\delta\leq\delta V_K^*+M.
\end{equation}

\begin{lemma}\label{lem:affine-growth-increment-bounds}
For $0<\alpha<\beta$, one has on $\mathbb C^n$
\begin{equation}\label{eq:affine-growth-increment-bounds}
 (\beta-\alpha)V_K^*\leq F_\beta-F_\alpha
 \leq\frac{\beta-\alpha}{\alpha}(F_\alpha-m).
\end{equation}
\end{lemma}
\begin{proof}
If $u$ is a competitor for $V_\alpha=V_{\alpha,\psi}$ and $v$ is a competitor for $V_K=V_{1,0}$, then $u+(\beta-\alpha)v$ is a competitor for $V_\beta$. Taking the suprema over $u$ and $v$ independently gives
\[
 V_\beta\geq V_\alpha+(\beta-\alpha)V_K.
\]
By \cref{prop:affine-projective-envelope}, the raw envelopes agree almost everywhere with their regularizations, so this inequality holds almost everywhere for the regularizations, and hence everywhere by \cref{lem:qpsh-almost-everywhere-order}. This is the first inequality. The same argument applied to \eqref{eq:affine-growth-raw-slope} gives the second.
\end{proof}

\begin{proposition}\label{prop:affine-envelope-slopes}
For every $a\in\mathbb C^n$ and $\tau>0$, the right derivative
\[
 D_\tau(a)=\lim_{h\downarrow0}\frac{F_{\tau+h}(a)-F_\tau(a)}{h}
\]
exists and is finite. The difference quotients decrease as $h>0$ increases, and
\begin{equation}\label{eq:affine-initial-slope-bounds}
 V_K^*(a)\leq\frac{F_{\tau+h}(a)-F_\tau(a)}h
 \leq D_\tau(a)\leq\frac{F_\tau(a)-m}{\tau}.
\end{equation}
As $h\to\infty$, they converge to $V_K^*(a)$, with the error bound
\begin{equation}\label{eq:affine-large-increment-error}
 0\leq\frac{F_{\tau+h}(a)-F_\tau(a)}h-V_K^*(a)
 \leq\frac{M-m}h.
\end{equation}
Moreover, $D_\tau(a)=V_K^*(a)$ if and only if
\begin{equation}\label{eq:affine-envelope-affine-ray}
 F_\delta(a)=F_\tau(a)+(\delta-\tau)V_K^*(a)
 \qquad(\delta\geq\tau).
\end{equation}
\end{proposition}
\begin{proof}
For $0<h_1<h_2$, concavity of $F_\delta(a)$ gives
\[
 F_{\tau+h_1}(a)-F_\tau(a)
 \geq\frac{h_1}{h_2}\bigl(F_{\tau+h_2}(a)-F_\tau(a)\bigr).
\]
Thus the difference quotients decrease with $h$. \Cref{lem:affine-growth-increment-bounds} bounds them between $V_K^*(a)$ and $(F_\tau(a)-m)/\tau$, which proves the existence and finiteness of $D_\tau(a)$ and \eqref{eq:affine-initial-slope-bounds}. Combining the lower bound for the increments with \eqref{eq:affine-obstacle-sandwich} gives \eqref{eq:affine-large-increment-error}, with a bound independent of $a$. If $D_\tau(a)=V_K^*(a)$, then by \eqref{eq:affine-initial-slope-bounds} every difference quotient equals $V_K^*(a)$, which gives \eqref{eq:affine-envelope-affine-ray}. Conversely, that formula gives $D_\tau(a)=V_K^*(a)$.
\end{proof}
\index{Envelope!growth derivative}

\begin{corollary}\label{cor:affine-slope-regular-contact}
Let $a\in K$ and $\tau\in\Lambda_{\mathrm{reg}}(a)$. Regular contact at $a$ holds for every $\delta>\tau$ if and only if $D_\tau(a)=0$. If $D_\tau(a)>0$, regular contact at $a$ fails for every $\delta>\tau$, and $\tau$ is the smallest element of $\Lambda$.
\end{corollary}
\begin{proof}
At $\tau$ one has $F_\tau(a)=\psi(a)$, and raw contact holds for all $\delta\geq\tau$ by \cref{thm:affine-growth-contact-ray}. If $D_\tau(a)=0$, then by \eqref{eq:affine-initial-slope-bounds} every difference quotient at $\tau$ vanishes, so $F_\delta(a)=\psi(a)$ for all $\delta>\tau$. If $D_\tau(a)>0$, then $F_\delta(a)>\psi(a)$ for all $\delta$ slightly larger than $\tau$, and by monotonicity for all $\delta>\tau$. In this case $\Lambda_{\mathrm{reg}}(a)=\{\tau\}$, so \cref{thm:affine-growth-contact-ray} shows that $\tau$ is the smallest element of $\Lambda$.
\end{proof}

Assume now that the weights are continuous. If the obstacle is imposed only quasi-everywhere on $K$, the regularized envelope itself becomes a competitor, so the supremum needs no regularization; a constraint imposed at every point would miss the values created by regularization. Indeed, for $q\in C(K,\mathbb R)$,
\begin{equation}\label{eq:affine-qe-pointwise-envelope}
 F_{\delta,q}(z)=\sup\{u(z):u\in\mathcal L_\delta,
                          \ u\leq q\text{ q.e. on }K\}.
\end{equation}
The envelope $F_{\delta,q}$ itself is a competitor on the right, since it agrees with the raw envelope off a pluripolar set by \cref{prop:affine-projective-envelope}. Conversely, the argument of \cref{cor:affine-polar-obstacle-invariance} shows that every competitor on the right lies below $F_{\delta,q}$ everywhere.

For $a\in\mathbb C^n$, let $\mathcal J_a(K)$ be the set of \emph{supporting pairs} $(\mu,c)$, where $\mu$ is a Borel probability measure on $K$ and $c\geq0$, such that
\begin{equation}\label{eq:affine-supporting-probability-tests}
 u(a)\leq\int_K f\,\mathrm{d}\mu+c
\end{equation}
whenever $u\in\mathcal L_1$, $f\in C(K,\mathbb R)$ and $u\leq f$ quasi-everywhere on $K$. We call $\mu$ a \emph{supporting probability} at $a$.

\begin{theorem}\label{thm:affine-supporting-probabilities}
For every $a\in\mathbb C^n$, $q\in C(K,\mathbb R)$ and $\delta>0$,
\begin{equation}\label{eq:affine-supporting-probability-minimum}
 F_{\delta,q}(a)=\min_{(\mu,c)\in\mathcal J_a(K)}
             \left(\int_K q\,\mathrm{d}\mu+\delta c\right).
\end{equation}
In particular, $\mathcal J_a(K)$ is nonempty and
\[
 \min_{(\mu,c)\in\mathcal J_a(K)}c=V_K^*(a).
\]
\end{theorem}
\begin{proof}
On the real Banach space $C(K)$, put $T(f)=F_{1,f}(a)$. By \eqref{eq:affine-qe-pointwise-envelope}, $T$ is concave, since convex combinations of competitors are competitors. Comparison of obstacles and translation by constants show that $T$ is nondecreasing and satisfies
\[
 T(f+t)=T(f)+t,\qquad |T(f)-T(g)|\leq\|f-g\|_\infty.
\]
Fix $f_0\in C(K)$. The directional derivative of the convex Lipschitz function $-T$ at $f_0$ is finite and sublinear. By the Hahn--Banach theorem, there is a continuous linear functional dominated by it, and since difference quotients of a convex function are monotone, this functional supports $-T$ at $f_0$. Changing the sign, we obtain a continuous linear functional $\ell$ with
\[
 T(f)\leq T(f_0)+\ell(f-f_0)\qquad(f\in C(K)).
\]
Testing $f=f_0+th$ with $h\geq0$ and $t>0$, and using monotonicity, gives $\ell(h)\geq0$. Testing with $f=f_0\pm t$ gives $\ell(1)=1$. By the Riesz representation theorem, $\ell(f)=\int_K f\,\mathrm{d}\mu$ for a Borel probability measure $\mu$ on $K$. Set $c=T(f_0)-\int_K f_0\,\mathrm{d}\mu$. The supporting inequality at $f=0$ gives $c\geq T(0)=V_K^*(a)\geq0$. Together with \eqref{eq:affine-qe-pointwise-envelope}, it shows that $(\mu,c)\in\mathcal J_a(K)$ and
\[
 T(f_0)=\int_K f_0\,\mathrm{d}\mu+c.
\]
Conversely, applying \eqref{eq:affine-supporting-probability-tests} to $u=F_{1,f_0}$, which lies below $f_0$ quasi-everywhere on $K$, shows that $T(f_0)\leq\int_Kf_0\,\mathrm{d}\mu+c$ for every pair in $\mathcal J_a(K)$. This proves \eqref{eq:affine-supporting-probability-minimum} for $\delta=1$, with the minimum attained. The identity $F_{\delta,q}=\delta F_{1,q/\delta}$ gives the general case, and the zero obstacle gives the last assertion.
\end{proof}
\index{Envelope!supporting probability}

\begin{theorem}\label{thm:affine-supported-growth-derivative}
If $\psi\in C(K,\mathbb R)$, then for every $a\in\mathbb C^n$ and $\tau>0$,
\begin{equation}\label{eq:affine-supported-growth-derivative}
 D_\tau(a)=\min\left\{c:(\mu,c)\in\mathcal J_a(K),
                \ \int_K\psi\,\mathrm{d}\mu+\tau c=F_\tau(a)\right\}.
\end{equation}
Consequently, $D_\tau(a)=V_K^*(a)$ if and only if there is a Borel probability measure $\mu$ on $K$ such that
\[
 (\mu,V_K^*(a))\in\mathcal J_a(K),\qquad
 \int_K\psi\,\mathrm{d}\mu=F_\tau(a)-\tau V_K^*(a).
\]
\end{theorem}
\begin{proof}
Each defining inequality of $\mathcal J_a(K)$ is preserved under weak convergence of the probability measures and convergence of the constants, because the test function $f$ is continuous. Since the Borel probability measures on the compact metric space $K$ form a weakly compact set, the pairs in $\mathcal J_a(K)$ with $c\leq C$ form a compact set for each finite $C$.

Call a pair in $\mathcal J_a(K)$ \emph{active} at $\tau$ if it satisfies the equality in \eqref{eq:affine-supported-growth-derivative}. Such pairs exist by \cref{thm:affine-supporting-probabilities}, and since $\int_K\psi\,\mathrm{d}\mu\geq m=\inf_K\psi$, they satisfy
\[
 0\leq c\leq\frac{F_\tau(a)-m}{\tau}.
\]
The set of active pairs is therefore compact, and the minimum of $c$ over it is attained. By \cref{thm:affine-supporting-probabilities} at $\tau+h$, each active pair gives $F_{\tau+h}(a)\leq F_\tau(a)+hc$, hence $D_\tau(a)\leq c$.

Conversely, for $h>0$ choose a pair $(\mu_h,c_h)$ attaining the minimum in \eqref{eq:affine-supporting-probability-minimum} at $\tau+h$, with $q=\psi$. The representation at $\tau$ gives
\[
 F_\tau(a)\leq\int_K\psi\,\mathrm{d}\mu_h+\tau c_h
            =F_{\tau+h}(a)-hc_h.
\]
Thus $0\leq c_h\leq(F_{\tau+h}(a)-F_\tau(a))/h$, and these constants are bounded by \eqref{eq:affine-initial-slope-bounds}. Along a sequence $h\downarrow0$, compactness gives a limiting pair $(\mu,c)\in\mathcal J_a(K)$. The same bound makes $F_{\tau+h}(a)\to F_\tau(a)$ and $hc_h\to0$, and continuity of $\psi$ then gives $\int_K\psi\,\mathrm{d}\mu+\tau c=F_\tau(a)$. Thus the limit pair is active, and it satisfies $c\leq D_\tau(a)$. This proves \eqref{eq:affine-supported-growth-derivative}. Finally, every pair has $c\geq V_K^*(a)$ by \cref{thm:affine-supporting-probabilities}, so the attainment of the minimum gives the stated criterion.
\end{proof}

At a point $a$ of regular contact, $F_\tau(a)=\psi(a)$, so the criterion involves only the original weight and tests by functions in $\mathcal L_1$. If the defining tests instead required $u\leq f$ at every point of $K$, the pair formed by the unit point mass at $a\in K$ and the constant $0$ would always be allowed. The resulting infimum would then be at most $\psi(a)$, so it could not detect a defect created by regularization.

\subsubsection{Variation of continuous weights}
\label{sec:affine-weight-response}

Supporting probabilities describe changes of the obstacle as well as the growth allowance. Keep the non-pluripolar compact $K\subset\mathbb C^n$ fixed. For $a\in\mathbb C^n$, $\delta>0$ and $q\in C(K,\mathbb R)$, put
\begin{equation}\label{eq:affine-active-supports}
 \mathcal A_a(\delta,q)=
 \left\{(\mu,c)\in\mathcal J_a(K):
   \int_Kq\,\mathrm{d}\mu+\delta c=F_{\delta,q}(a)\right\}.
\end{equation}
The weights carry the uniform norm, probabilities carry the weak topology, and constants carry the ordinary topology.

\begin{lemma}\label{lem:affine-active-support-stability}
The set $\mathcal A_a(\delta,q)$ is nonempty, convex and compact. For $\delta,\eta\geq\epsilon>0$ and $p,q\in C(K,\mathbb R)$,
\begin{equation}\label{eq:affine-weight-growth-lipschitz}
 |F_{\eta,p}(a)-F_{\delta,q}(a)|
 \leq\|p-q\|_\infty+
 |\eta-\delta|\left(V_K^*(a)+
                    \frac{\operatorname{osc}_Kq}{\epsilon}\right).
\end{equation}
If $\delta_j\to\delta>0$, $q_j\to q$ uniformly and $(\mu_j,c_j)\in\mathcal A_a(\delta_j,q_j)$, then these pairs have a convergent subsequence, and every subsequential limit belongs to $\mathcal A_a(\delta,q)$.
\end{lemma}
\begin{proof}
Nonemptiness follows from \cref{thm:affine-supporting-probabilities}. The continuous test inequalities make $\mathcal J_a(K)$ closed and convex. The active equality is also closed and affine. For an active pair, \eqref{eq:affine-obstacle-sandwich} gives
\[
 0\leq c\leq\frac{F_{\delta,q}(a)-\inf_Kq}{\delta}
 \leq V_K^*(a)+\frac{\operatorname{osc}_Kq}{\delta}.
\]
The bounded constant slice is compact, as in the proof of \cref{thm:affine-supported-growth-derivative}; hence the active set is compact.

With $q$ fixed, \cref{lem:affine-growth-increment-bounds} and \eqref{eq:affine-obstacle-sandwich} bound a growth increment by its length times $V_K^*(a)+\operatorname{osc}_Kq/\epsilon$. Obstacle order and constant translation bound the change from $q$ to $p$ at fixed $\eta$ by $\|p-q\|_\infty$. Their sum is \eqref{eq:affine-weight-growth-lipschitz}.

For the sequence in the statement, the active constants have a common bound because the weights are uniformly bounded and the growth allowances stay bounded below. Weak compactness of probabilities gives subsequential pair convergence. Closedness retains membership in $\mathcal J_a(K)$. The displayed Lipschitz estimate gives envelope-value convergence, while uniform convergence of $q_j$ and weak convergence of $\mu_j$ give convergence of the integrals. Passing to the active equality proves the last assertion.
\end{proof}
\index{Envelope!weight variation}

\begin{proposition}\label{prop:affine-joint-directional-response}
For $r\in C(K,\mathbb R)$ and $s\in\mathbb R$, one has
\begin{equation}\label{eq:affine-joint-directional-response}
 \lim_{h\downarrow0}
 \frac{F_{\delta+hs,q+hr}(a)-F_{\delta,q}(a)}h
 =\min_{(\mu,c)\in\mathcal A_a(\delta,q)}
       \left(\int_Kr\,\mathrm{d}\mu+sc\right).
\end{equation}
The same limit holds along $h_j\downarrow0$ with $s_j\to s$ and $r_j\to r$ uniformly, replacing $s,r$ in the quotient by $s_j,r_j$. Thus the one-sided directional derivative is stable under convergent directions.
\end{proposition}
\begin{proof}
Compactness of the active set attains the minimum on the right. For a fixed base active pair, the representation at the perturbed parameters gives
\[
 F_{\delta+hs,q+hr}(a)
 \leq F_{\delta,q}(a)+h\left(\int_Kr\,\mathrm{d}\mu+sc\right).
\]
A minimizing base pair therefore bounds the limsup by the stated minimum. Conversely, choose a minimizing pair $(\mu_h,c_h)$ at $(\delta+hs,q+hr)$. The representation at the base parameters gives
\[
 \frac{F_{\delta+hs,q+hr}(a)-F_{\delta,q}(a)}h
 \geq\int_Kr\,\mathrm{d}\mu_h+sc_h.
\]
The local Lipschitz estimate bounds the quotients. Along a sequence approaching their liminf, \cref{lem:affine-active-support-stability} supplies a convergent subsequence of pairs with a base active limit. Its directional value is at least the minimum. This proves \eqref{eq:affine-joint-directional-response}.

For convergent directions, the same two inequalities hold. The fixed minimizing base pair gives a convergent upper bound. Perturbed minimizing pairs still have uniformly bounded constants and base active limits. Uniform convergence of $r_j$ controls their moving integrals, and $s_j\to s$ controls the constant terms. The lower bound therefore has the same limit. All perturbed growth allowances are positive for sufficiently small increments.
\end{proof}

\begin{corollary}\label{cor:affine-two-sided-growth-slopes}
For fixed $a$ and $q$, the function $\delta\mapsto F_{\delta,q}(a)$ has finite left and right derivatives $D^-_{\delta,q}(a)$ and $D^+_{\delta,q}(a)$ at every positive parameter, with
\begin{equation}\label{eq:affine-two-sided-growth-slopes}
 D^-_{\delta,q}(a)=\max_{\mathcal A_a(\delta,q)}c,
 \qquad
 D^+_{\delta,q}(a)=\min_{\mathcal A_a(\delta,q)}c.
\end{equation}
It is differentiable exactly when all active pairs have the same constant $c$. For each fixed $a,q$, the set of positive parameters where it fails to be differentiable is at most countable.
\end{corollary}
\begin{proof}
Apply \eqref{eq:affine-joint-directional-response} with $r=0$ and $s=1,-1$. The negative direction gives minus the left derivative, proving \eqref{eq:affine-two-sided-growth-slopes}. Concavity gives
\[
 D^+_{\alpha,q}(a)\geq D^-_{\beta,q}(a)
 \qquad(0<\alpha<\beta).
\]
Hence the nonempty open intervals $(D^+_{\delta,q}(a),D^-_{\delta,q}(a))$ at exceptional parameters are pairwise disjoint. Assigning a rational number in each interval proves countability.
\end{proof}
\index{Envelope!growth differentiability}

\begin{corollary}\label{cor:affine-weight-differentiability}
At fixed $a$ and $\delta>0$, the map $q\mapsto F_{\delta,q}(a)$ on $C(K,\mathbb R)$ is G\^ateaux differentiable at $q$ if and only if the active probability is unique. In that case its derivative is
\[
 r\longmapsto\int_Kr\,\mathrm{d}\mu.
\]
\end{corollary}
\begin{proof}
With $s=0$, \eqref{eq:affine-joint-directional-response} gives the minimum of $\int_Kr\,\mathrm{d}\mu$ over active pairs. The negative direction gives minus the maximum. If the active probability is unique, these values give the same continuous linear two-sided derivative. Conversely, a linear two-sided derivative forces the minimum and maximum to agree for every continuous $r$. Such tests separate Borel probabilities on $K$, so all active probabilities coincide. Their constant is then fixed by \eqref{eq:affine-active-supports} and $\delta>0$.
\end{proof}
\index{Envelope!weight differentiability}

\begin{example}\label{ex:affine-quadratic-weight-corner}
Let $K=\overline D\subset\mathbb C$ be the closed unit disk and fix $\delta>0$. For $q_t(z)=-t|z|^2$,
\begin{equation}\label{eq:affine-quadratic-weight-corner}
 F_{\delta,q_t}(0)=
 \begin{cases}
 -t,&t\geq0,\\
 0,&-\delta/2\leq t\leq0.
 \end{cases}
\end{equation}
Thus the weight functional at the origin is not G\^ateaux differentiable at the zero weight, and it has more than one active probability there.
\end{example}
\begin{proof}
For $t\geq0$, each raw competitor is at most $-t$ on the boundary circle. The subharmonic maximum principle bounds it by $-t$ on $D$. Regularization retains that bound near zero, and the constant $-t$ is a global competitor.

For $t=-b$ with $0\leq b\leq\delta/2$, define
\[
 u_b(z)=
 \begin{cases}
 b|z|^2,&|z|\leq1,\\
 b+\delta\log|z|,&|z|>1.
 \end{cases}
\]
This function is continuous and subharmonic. Its ordinary Laplacian has interior density $4b$, zero exterior density and a nonnegative circle contribution from the outward derivative jump $\delta-2b$. It belongs to $\mathcal L_\delta$ and meets $q_{-b}$ on $K$. Hence $F_{\delta,q_{-b}}(0)\geq0$. On the open disk the raw envelope is at most the continuous function $b|z|^2$, so regularization gives the reverse inequality at zero. This proves \eqref{eq:affine-quadratic-weight-corner}. The right and left derivatives in $t$ are $-1$ and $0$. \Cref{cor:affine-weight-differentiability} gives the final assertion.
\end{proof}


\subsubsection{Local regularity and strict weights}
\label{sec:affine-local-regularity}

Local contact depends on the compact pieces near the point. The relative extremal function of \cref{def:relative-extremal-function} detects their unweighted regularity. A strict lower Hessian bound on a global extending weight then turns failure of that regularity into a quantitative weighted defect.

For a non-pluripolar compact $E\subset\mathbb C^n$, write $V_E$ and $V_E^*$ for its raw and regularized growth-one envelopes with zero obstacle. For other weights and growth allowances, write $F_{\delta,E,q}$ for the regularized envelope on $E$ with the everywhere obstacle $q$.

\begin{lemma}\label{lem:relative-affine-zero-value}
Let $0<r<R$, $B=B(a,R)$, and let $E\subseteq\overline B(a,r)$ be compact. If $E$ is pluripolar, then $h_{E,B}^*=0$ on $B$. If $E$ is non-pluripolar, then
\[
 h_{E,B}^*(a)=-1\quad\Longleftrightarrow\quad V_E^*(a)=0.
\]
\end{lemma}
\begin{proof}
In the pluripolar case, choose $b\in\mathcal L_1$ with poles on $E$ by \cref{lem:affine-logarithmic-pole-barrier}, and choose a finite upper bound $C$ for $b$ on $B$. The relative competitors
\[
 u_\epsilon=\max\{-1,\epsilon(b-C)\}
\]
tend to zero almost everywhere as $\epsilon\downarrow0$. Their upper bound $h_{E,B}^*$ is therefore nonnegative almost everywhere, then everywhere by canonical order. Since it is also nonpositive, it vanishes.

Suppose $E$ is non-pluripolar. \Cref{prop:affine-projective-envelope} gives a finite upper bound $M\geq1$ for $V_E^*$ on $B$. For each raw global competitor $v$, the function $\max(-1,v/M-1)$ is a relative competitor. Taking suprema and regularizing gives
\[
 h_{E,B}^*(a)\geq\max\{-1,V_E^*(a)/M-1\}.
\]
As $V_E^*\geq0$, the relative value $-1$ forces the global value zero.

For the converse choose $r<R_0<R$ and $k>0$ with $k\log(R_0/r)>1$. The plurisubharmonic function
\[
 L(z)=k\log\bigl(\lVert z-a\rVert/r\bigr)
\]
has growth allowance $k$, with value $-\infty$ at $a$. For a relative competitor $u$, paste $W=\max(u+1,L)$ on $B$ with $W=L$ outside $B$. On the outer collar $L>1\geq u+1$, so the pieces agree and \cref{lem:qpsh-gluing} applies. On $E$ one has $W=0$; on the open inner ball one has $W=u+1$. Hence $W/k$ is a global growth-one competitor, and the same $k$ works for the whole relative family. Supremum and regularization near $a$ give
\[
 V_E^*(a)\geq(1+h_{E,B}^*(a))/k.
\]
A zero global value thus forces the relative value $-1$.
\end{proof}

Fix again a non-pluripolar compact $K\subset\mathbb C^n$. For $a\in K$ put $K_r=K\cap\overline B(a,r)$.

\begin{definition}\label{def:affine-local-pluriregularity}\index{Pluriregular compact!local}
The compact $K$ is locally pluriregular at $a$ if every $K_r$, $r>0$, is non-pluripolar and $V_{K_r}^*(a)=0$. It is locally $(\delta,q)$-regular there if $F_{\delta,K_r,q}(a)=q(a)$ for every $r>0$.
\end{definition}

It suffices to check arbitrarily small radii: compact-set monotonicity handles the larger pieces. \Cref{lem:relative-affine-zero-value} identifies local pluriregularity with the relative values $h_{K_r,B(a,R)}^*(a)=-1$ on small concentric balls.

\begin{lemma}\label{lem:affine-local-weighted-regularity}
Let $q$ be bounded and upper semicontinuous on $K$, and suppose $\delta>0$ is a global raw-contact parameter for $q$. If $K$ is locally pluriregular at $a$, then $F_{\delta,K,q}(a)=q(a)$ and $K$ is locally $(\delta,q)$-regular there. These equalities persist at every later growth parameter.
\end{lemma}
\begin{proof}
Put $M_r=\sup_{K_r}q$. Raw contact, compact-set monotonicity and the obstacle sandwich give
\[
 q(a)\leq F_{\delta,K,q}(a)
 \leq F_{\delta,K_r,q}(a)
 \leq M_r+\delta V_{K_r}^*(a)=M_r.
\]
Since $q$ is upper semicontinuous and $a\in K_r$, one has $M_r\downarrow q(a)$, giving the global equality.

Every compact subset $E\subseteq K$ inherits raw contact: on $E$ its raw envelope lies between the original envelope $q$ and its obstacle $q$. For fixed $s>0$, apply the same bounds with $r<s$ and $K_s$ in place of $K$, then let $r\downarrow0$. This proves local weighted regularity. The raw contact ray contains every later parameter by \cref{thm:affine-growth-contact-ray}; applying the proved criterion at each of them gives persistence.
\end{proof}

\begin{theorem}\label{thm:affine-local-strict-weight}
Let $a\in K$, $\tau>0$ and $\Psi\in\mathcal L_\tau\cap C(\mathbb C^n)$, and put $q=\Psi|_K$. Suppose that, for some $R,c>0$,
\[
 \ddc\Psi\geq c\omega_0\quad\text{on }B(a,R),
 \qquad\omega_0=\ddc\lVert z\rVert^2.
\]
For $0<r<R$ and $\delta\geq\tau$,
\begin{equation}\label{eq:affine-local-strict-weight-defect}
 F_{\delta,K,q}(a)-\Psi(a)
 \geq cr^2\bigl(1+h_{K_r,B(a,R)}^*(a)\bigr).
\end{equation}
Regular contact at $a$ for one $\delta\geq\tau$ is equivalent to local pluriregularity at $a$. In that case regular contact holds for every $\delta\geq\tau$.
\end{theorem}
\begin{proof}
Set $B=B(a,R)$ and $E=K_r$. For a relative competitor $u$, the function
\[
 H_u(z)=\Psi(z)-c\lVert z-a\rVert^2+cr^2(u(z)+1)
\]
is plurisubharmonic on $B$. Indeed, the current inequality makes its first two terms plurisubharmonic, and the last term is plurisubharmonic. Since $u\leq0$, one has $H_u\leq\Psi$ where $\lVert z-a\rVert\geq r$. Paste $W=\max(\Psi,H_u)$ on $B$ with $W=\Psi$ outside $B$. The pieces agree on an outer collar, so $W$ is globally plurisubharmonic. It belongs to $\mathcal L_\tau$ because it equals $\Psi$ outside the ball. On $E$, $u=-1$ gives $H_u\leq\Psi$; on $K\setminus E$ the distance bound gives the same inequality. Thus $W=q$ everywhere on $K$ and is a raw competitor for each $\delta\geq\tau$. Taking the supremum of $H_u$ near $a$ and regularizing, with the continuous $\Psi$ retained, proves \eqref{eq:affine-local-strict-weight-defect}.

The global extension $\Psi$ itself gives raw contact for every $\delta\geq\tau$. If regular contact holds at one such parameter, the left side of the estimate is zero. Each relative value must then be $-1$. \Cref{lem:relative-affine-zero-value} excludes polar pieces and gives $V_{K_r}^*(a)=0$ for all $r<R$; larger pieces follow by monotonicity. Conversely, local pluriregularity and \cref{lem:affine-local-weighted-regularity} give regular contact at every $\delta\geq\tau$.
\end{proof}
\index{Envelope!locally strict weight}

In this regular case, $V_K^*(a)=0$ by compact-set monotonicity and the initial growth derivative at $\tau$ is zero by persistence and \cref{cor:affine-slope-regular-contact}.

\Needspace{10\baselineskip}
\begin{corollary}\label{cor:affine-weight-degenerate-hessian}\index{Contact set!weight Hessian degeneracy}
Let $\tau>0$, $\Psi\in\mathcal L_\tau\cap C^2(\mathbb C^n)$ and $a\in K$. If $F_{\tau,K,\Psi|_K}(a)=\Psi(a)$ but regular contact fails at some later parameter, then
\[
 (\ddc\Psi)^n(a)=0.
\]
\end{corollary}
\begin{proof}
The smooth Hessian of $\Psi$ is semipositive. A nonzero top power at $a$ would make it positive definite. Continuity would then give $\ddc\Psi\geq c\omega_0$ on a neighborhood of $a$ for some $c>0$. \Cref{thm:affine-local-strict-weight} would force persistence of regular contact, a contradiction.
\end{proof}




\subsubsection{A regularization defect along a complex disk}
\label{sec:affine-contact-defect}
\index{Contact set!regularization defect}

Raw contact with a constant obstacle can hold for every positive growth parameter while regularization raises the value of the envelope at a point of the compact set. In the following example, a complex disk through the point, transverse to the planar direction, prevents any local plurisubharmonic function from separating the point from the rest of the compact set, but it does not remove this defect.

In the complex plane, let $D$ be the open unit disk and put
\[
 p_j=2+2^{-j},\qquad r_j=\exp(-16\,2^j),\qquad
 E=\overline D\cup\{2\}\cup\bigcup_{j\geq1}\overline D(p_j,r_j).
\]

\Needspace{17\baselineskip}
\begin{lemma}\label{lem:planar-contact-barrier}
The set $E$ is compact, and every neighborhood of each of its points contains an open disk contained in $E$. The series
\[
 h(\zeta)=\sum_{j\geq1}2^{-j}\log|\zeta-p_j|+2\log2
\]
defines a subharmonic function on $\mathbb C$ of logarithmic growth at most one; it satisfies $h(2)=0$ and is at most $-16+2\log2$ on every small disk $\overline D(p_j,r_j)$. The function
\[
 v(\zeta)=\frac{8\log^+|\zeta|+h(\zeta)-4\log2}{9}
\]
is subharmonic of logarithmic growth at most one, is nonpositive on $E\setminus\{2\}$, and satisfies $v(2)=4\log2/9$.
\end{lemma}
\begin{proof}
Since $r_j<2^{-j-3}$, the small disks are pairwise disjoint and disjoint from $\overline D$, and they accumulate only at $2$. This proves compactness. Every neighborhood of $2$ contains one of the small disks, and every other point of $E$ lies in the closure of the interior of one of the disks.

Fix $R>0$. On the disk $\{|\zeta|<R\}$ every term satisfies $\log|\zeta-p_j|<\log(R+3)$, so the partial sums of $\sum_j2^{-j}(\log|\zeta-p_j|-\log(R+3))$ decrease there. Their limit is therefore subharmonic or identically $-\infty$ on this disk, and it is finite at $\zeta=0$. Since $\sum_j2^{-j}=1$, this limit is $h-2\log2-\log(R+3)$, so $h$ is subharmonic on $\mathbb C$. The bound $|\zeta-p_j|\leq|\zeta|+3$ gives $h(\zeta)\leq\log(|\zeta|+3)+2\log2$, hence the growth bound, and $\sum_jj2^{-j}=2$ gives $h(2)=0$.

On the $j$th small disk, all distances to the other centers are less than one, so the corresponding terms are negative, and
\[
 h\leq2^{-j}\log r_j+2\log2=-16+2\log2.
\]
On $\overline D$ all distances to the centers are less than four, hence $h\leq4\log2$ and $v\leq0$. On each small disk $|\zeta|<4$, so $9v\leq14\log2-16<0$. The growth bound for $v$ and its value at $2$ follow directly from the formula.
\end{proof}

Now let $n\geq2$ and $K=E\times\overline D^{\,n-1}\subset\mathbb C^n$. For a non-pluripolar compact set $L\subset\mathbb C^n$, write $V_{\delta,L}$ and $F_{\delta,L}$ for the raw envelope with growth allowance $\delta$ and zero obstacle on $L$, and for its regularization; thus $F_{\delta,L}=F_{\delta,L,0}$ in the notation of \cref{sec:affine-local-regularity}.

\Needspace{19\baselineskip}
\begin{proposition}\label{prop:affine-transverse-contact-defect}
For every point of $K$ and every neighborhood $W$ of it, the set $K\cap W$ is non-pluripolar. For $\delta>0$ and $a=(2,w)$ with $w\in\overline D^{\,n-1}$,
\begin{equation}\label{eq:affine-contact-defect}
 V_{\delta,K}|_K=0,\qquad
 F_{\delta,K}(a)\geq\frac{4\delta}{9}\log2>0.
\end{equation}
For every $r>0$, the same lower bound holds for $F_{\delta,K\cap\overline B(a,r)}(a)$. At $a_0=(2,0,\ldots,0)$, every plurisubharmonic function $f$ defined near $a_0$ satisfies
\begin{equation}\label{eq:transverse-psh-no-jump}
 f(a_0)\leq\limsup_{K\setminus\{a_0\}\ni z\to a_0}f(z).
\end{equation}
Thus no plurisubharmonic function near $a_0$ separates $a_0$ from $K\setminus\{a_0\}$, that is, $K$ is not plurithin at $a_0$. Nevertheless, the regularized envelope does not touch the obstacle at $a_0$ for any positive growth parameter.
\end{proposition}
\index{Envelope!non-thin contact defect}
\begin{proof}
By the preceding lemma, every neighborhood of a point of $K$ contains a nonempty open subset of $K$, namely a product of open disks. This proves the first assertion, which applies in particular to each $K\cap\overline B(a,r)$, $r>0$.

The plurisubharmonic function $z\mapsto v(z_1)$ belongs to $\mathcal L_1$ and is nonpositive on $K$ outside the pluripolar set $K\cap\{z_1=2\}$. By \cref{cor:affine-polar-obstacle-invariance} with zero obstacle, $v(z_1)\leq F_{1,K}(z)$ for all $z$. The value $v(2)=4\log2/9$ gives the lower bound for $\delta=1$. Since $0$ is a competitor, $V_{\delta,K}|_K=0$, and scaling of competitors gives $F_{\delta,K}=\delta F_{1,K}$. Replacing $K$ by a smaller compact set only increases the envelope, which proves the local assertion.

Finally, let $f$ be plurisubharmonic near $a_0$, and restrict it to the complex line through $a_0$ in which only the second coordinate varies. Small circles in this line centered at $a_0$ lie in $K\setminus\{a_0\}$, and the submean inequality on these circles gives \eqref{eq:transverse-psh-no-jump}.
\end{proof}

